% Options for packages loaded elsewhere
\PassOptionsToPackage{unicode}{hyperref}
\PassOptionsToPackage{hyphens}{url}
\documentclass[
]{article}
\usepackage{xcolor}
\usepackage{amsmath,amssymb}

\usepackage{hyperref}
\usepackage{ulem}  % For underlining links
\normalem

\hypersetup{
    colorlinks=false,         % Disable colored links (boxes will appear)
    citebordercolor=red,      % Border color for citation links
    linkbordercolor=red,      % Border color for internal links
    urlbordercolor=blue,      % Border color for external links (URLs)
    pdfborder={0 0 2}         % Specifies the border thickness: {horizontal vertical thickness}
}


% Only underline external links (URLs)
\let\oldhref\href
\renewcommand{\href}[2]{\ifx#1\urlprefix\oldhref{#1}{#2}\else\uline{\oldhref{#1}{#2}}\fi}


\renewcommand{\[}{\begin{equation}}
\renewcommand{\]}{\end{equation}}


\setcounter{secnumdepth}{5}
\usepackage{iftex}
\ifPDFTeX
  \usepackage[T1]{fontenc}
  \usepackage[utf8]{inputenc}
  \usepackage{textcomp} % provide euro and other symbols
\else % if luatex or xetex
  \usepackage{unicode-math} % this also loads fontspec
  \defaultfontfeatures{Scale=MatchLowercase}
  \defaultfontfeatures[\rmfamily]{Ligatures=TeX,Scale=1}
\fi
\usepackage{lmodern}
\ifPDFTeX\else
  % xetex/luatex font selection
\fi
% Use upquote if available, for straight quotes in verbatim environments
\IfFileExists{upquote.sty}{\usepackage{upquote}}{}
\IfFileExists{microtype.sty}{% use microtype if available
  \usepackage[]{microtype}
  \UseMicrotypeSet[protrusion]{basicmath} % disable protrusion for tt fonts
}{}
\makeatletter
\@ifundefined{KOMAClassName}{% if non-KOMA class
  \IfFileExists{parskip.sty}{%
    \usepackage{parskip}
  }{% else
    \setlength{\parindent}{0pt}
    \setlength{\parskip}{6pt plus 2pt minus 1pt}}
}{% if KOMA class
  \KOMAoptions{parskip=half}}
\makeatother
\setlength{\emergencystretch}{3em} % prevent overfull lines
\providecommand{\tightlist}{%
  \setlength{\itemsep}{0pt}\setlength{\parskip}{0pt}}
\usepackage[]{biblatex}
\addbibresource{refs.bib}
\usepackage{bookmark}
\IfFileExists{xurl.sty}{\usepackage{xurl}}{} % add URL line breaks if available
\urlstyle{same}

\title{Pion Dressing and Magnetic Transitions}
  \author{Ida}
  \date{\today}  % Default to today if no date is provided

\begin{document}
\maketitle

\section{Toy Model of Pion Dressing and Magnetic
Transitions}\label{toy-model-of-pion-dressing-and-magnetic-transitions}

These notes isolate one specific piece of physics:

\begin{enumerate}
\def\labelenumi{\arabic{enumi}.}
\tightlist
\item
  start from a toy Hamiltonian with explicit pion and photon modes,
\item
  integrate out the pionful sector,
\item
  diagonalise the resulting low-energy nuclear Hamiltonian,
\item
  show that the quantised photon still couples the resulting
  pion-dressed nuclear states.
\end{enumerate}

The goal is pedagogical rather than realistic. We keep the strong
dynamics schematic, but we keep the logic of the reduction explicit. In
particular, we do \textbf{not} yet try to identify the states with
realistic deuteron channels. That can be added later once the toy
mechanism is clear.

Throughout, we set \(\hbar = 1\).

\subsection{1) Degrees of freedom and physical
picture}\label{degrees-of-freedom-and-physical-picture}

We introduce three nuclear basis states:

\[
|g\rangle, \qquad |e\rangle, \qquad |\chi\rangle.
\]

The intended interpretation is:

\begin{itemize}
\tightlist
\item
  \(|g\rangle\) and \(|e\rangle\) are two low-energy \textbf{no-pion}
  nuclear basis states that we want to keep.
\item
  \(|\chi\rangle\) is a higher-energy \textbf{doorway configuration}
  that will only appear together with one explicit pion.
\end{itemize}

We also introduce two bosonic modes:

\[
a, a^\dagger \qquad \text{for the photon mode,}
\]

\[
b, b^\dagger \qquad \text{for the pion mode.}
\]

The basis states of the full toy problem therefore look like

\[
|g,n_\gamma,0_\pi\rangle,\qquad |e,n_\gamma,0_\pi\rangle,\qquad |\chi,n_\gamma,1_\pi\rangle.
\]

The important point is that \(|g\rangle\) and \(|e\rangle\) are
\textbf{not yet} the physical ground and excited states. They are only
basis states in the retained low-energy sector. The physical low-energy
states will appear only after the pionful sector has been removed and
the effective \(2\times 2\) Hamiltonian has been diagonalised.

\subsection{2) Full toy Hamiltonian}\label{full-toy-hamiltonian}

We split the Hamiltonian into a diagonal part plus pion and photon
couplings:

\[
H = H_0 + V_{\pi N} + V_{\gamma N}.
\label{eq:Hfullsplit}
\]

The diagonal part is

\[
H_0
=
E_g |g\rangle\langle g|
+
E_e |e\rangle\langle e|
+
E_\chi |\chi\rangle\langle \chi|
+
\omega_\pi b^\dagger b
+
\omega_\gamma a^\dagger a.
\label{eq:H0toy}
\]

The pion coupling is chosen to connect the two low-energy basis states
to the same pionful doorway state:

\[
V_{\pi N}
=
\Big(
W_g |\chi\rangle\langle g|
+
W_e |\chi\rangle\langle e|
\Big)b^\dagger
+
b\Big(
W_g |g\rangle\langle \chi|
+
W_e |e\rangle\langle \chi|
\Big).
\label{eq:Vpiontoy}
\]

Thus \(|g,0_\pi\rangle\) and \(|e,0_\pi\rangle\) can both make virtual
excursions into \(|\chi,1_\pi\rangle\).

The photon couples through a magnetic operator \(\hat\mu\):

\[
V_{\gamma N}
=
g_\gamma (a+a^\dagger)\,\hat\mu.
\label{eq:Vphotonfull}
\]

For the main derivation it is enough to keep the magnetic operator
inside the retained low-energy space:

\[
\hat\mu
=
\mu_g |g\rangle\langle g|
+
\mu_e |e\rangle\langle e|
+
m_0\Big(|g\rangle\langle e| + |e\rangle\langle g|\Big).
\label{eq:muBare}
\]

Two special cases are worth keeping in mind:

\begin{itemize}
\tightlist
\item
  \(m_0 = 0\): the bare operator is diagonal in the
  \(\{|g\rangle,|e\rangle\}\) basis.
\item
  \(m_0 \neq 0\): there is already a direct magnetic transition in that
  basis.
\end{itemize}

The first case is the cleanest one for seeing how pion-induced mixing
alone can generate a transition between the dressed states.

\subsection{3) Retained and eliminated
sectors}\label{retained-and-eliminated-sectors}

We keep the no-pion sector and eliminate the one-pion doorway sector.

The retained projector is

\[
P
=
\sum_{n_\gamma=0}^{\infty}
\Big(
|g,n_\gamma,0_\pi\rangle\langle g,n_\gamma,0_\pi|
+
|e,n_\gamma,0_\pi\rangle\langle e,n_\gamma,0_\pi|
\Big),
\label{eq:Pprojector}
\]

and the eliminated projector is

\[
Q
=
\sum_{n_\gamma=0}^{\infty}
|\chi,n_\gamma,1_\pi\rangle\langle \chi,n_\gamma,1_\pi|.
\label{eq:Qprojector}
\]

The exact projected Schrödinger equation is

\[
H_{\rm eff}(E)\,|\psi_P\rangle = E\,|\psi_P\rangle,
\label{eq:BHeigen}
\]

with the exact energy-dependent effective Hamiltonian

\[
H_{\rm eff}(E)
=
PHP + PHQ\,\frac{1}{E-QHQ}\,QHP.
\label{eq:BHtoy}
\]

Eq. \(\ref{eq:BHtoy}\) is exact. The approximation will enter only when
we expand it to low order in the pion coupling.

\subsection{4) Second-order elimination of the pionful
sector}\label{second-order-elimination-of-the-pionful-sector}

We now keep terms through second order in \(W_g\) and \(W_e\).

The intermediate pionful energies seen from the two low-energy basis
states are

\[
\Delta_g = E_\chi + \omega_\pi - E_g,
\qquad
\Delta_e = E_\chi + \omega_\pi - E_e.
\label{eq:gapDefs}
\]

For low-energy states one has \(\Delta_g > 0\) and \(\Delta_e > 0\).

At second order, it is consistent to replace the exact eigenvalue \(E\)
in the denominator of Eq. \(\ref{eq:BHtoy}\) by an unperturbed
low-energy energy. The reason is simple: the resolvent term is already
\(O(W^2)\), while the difference between \(E\) and the unperturbed
energy is itself \(O(W^2)\), so feeding that correction back into the
denominator would only change the answer at \(O(W^4)\).

With that approximation, the pionful sector produces:

\begin{enumerate}
\def\labelenumi{\arabic{enumi}.}
\tightlist
\item
  a shift of the \(|g\rangle\) energy through the virtual process
\end{enumerate}

\[
|g,0_\pi\rangle \to |\chi,1_\pi\rangle \to |g,0_\pi\rangle,
\]

\begin{enumerate}
\def\labelenumi{\arabic{enumi}.}
\setcounter{enumi}{1}
\tightlist
\item
  a shift of the \(|e\rangle\) energy through
\end{enumerate}

\[
|e,0_\pi\rangle \to |\chi,1_\pi\rangle \to |e,0_\pi\rangle,
\]

\begin{enumerate}
\def\labelenumi{\arabic{enumi}.}
\setcounter{enumi}{2}
\tightlist
\item
  an induced mixing through
\end{enumerate}

\[
|g,0_\pi\rangle \to |\chi,1_\pi\rangle \to |e,0_\pi\rangle.
\]

The resulting effective nuclear Hamiltonian in the retained
\(\{|g\rangle,|e\rangle\}\) basis is

\[
H_{\rm nuc}^{\rm eff}
=
\begin{pmatrix}
E_g - \dfrac{W_g^2}{\Delta_g}
&
-\dfrac{W_gW_e}{2}
\left(
\dfrac{1}{\Delta_g}
+
\dfrac{1}{\Delta_e}
\right)
\\[1.2em]
-\dfrac{W_gW_e}{2}
\left(
\dfrac{1}{\Delta_g}
+
\dfrac{1}{\Delta_e}
\right)
&
E_e - \dfrac{W_e^2}{\Delta_e}
\end{pmatrix}.
\label{eq:HeffNuclear}
\]

It is convenient to define

\[
E_g' = E_g - \frac{W_g^2}{\Delta_g},
\qquad
E_e' = E_e - \frac{W_e^2}{\Delta_e},
\label{eq:dressedDiagE}
\]

and

\[
V_\pi
=
-\frac{W_gW_e}{2}
\left(
\frac{1}{\Delta_g}
+
\frac{1}{\Delta_e}
\right).
\label{eq:VpiEff}
\]

Then Eq. \(\ref{eq:HeffNuclear}\) becomes

\[
H_{\rm nuc}^{\rm eff}
=
\begin{pmatrix}
E_g' & V_\pi \\
V_\pi & E_e'
\end{pmatrix}.
\label{eq:HeffCompact}
\]

This is the central structural result. After the pion has been
integrated out, there is no explicit pion basis vector left, but there
is now a direct low-energy coupling \(V_\pi\) between the retained
states.

\subsection{5) Pion-dressed nuclear
eigenstates}\label{pion-dressed-nuclear-eigenstates}

The physical low-energy nuclear states are obtained by diagonalising Eq.
\(\ref{eq:HeffCompact}\).

Write

\[
|1\rangle
=
\cos\theta\,|g\rangle
-
\sin\theta\,|e\rangle,
\qquad
|2\rangle
=
\sin\theta\,|g\rangle
+
\cos\theta\,|e\rangle.
\label{eq:dressedStates}
\]

The mixing angle is defined by

\[
\tan 2\theta = \frac{2V_\pi}{E_e' - E_g'}.
\label{eq:tan2theta}
\]

The corresponding dressed energies are

\[
E_{1,2}
=
\frac{E_g' + E_e'}{2}
\mp
\sqrt{
\left(
\frac{E_e' - E_g'}{2}
\right)^2
+
V_\pi^2
}.
\label{eq:dressedEnergies}
\]

At this point the interpretation changes:

\begin{itemize}
\tightlist
\item
  \(|g\rangle\) and \(|e\rangle\) are no longer the observable
  low-energy levels.
\item
  \(|1\rangle\) and \(|2\rangle\) are the pion-dressed low-energy
  states.
\end{itemize}

This is the toy version of the statement that the explicit pionful
dynamics has been absorbed into an effective no-pion Hamiltonian.

\subsection{6) Magnetic operator in the dressed
basis}\label{magnetic-operator-in-the-dressed-basis}

We now ask how the photon couples the dressed states.

In the retained basis, the bare magnetic operator was

\[
\hat\mu
=
\begin{pmatrix}
\mu_g & m_0 \\
m_0 & \mu_e
\end{pmatrix}_{\{|g\rangle,|e\rangle\}}.
\label{eq:muBareMatrix}
\]

The same rotation that diagonalises Eq. \(\ref{eq:HeffCompact}\)
transforms the magnetic operator to the dressed basis:

\[
\hat\mu_{\rm dressed} = U^\dagger \hat\mu U,
\label{eq:muDressedDef}
\]

where \(U\) is the rotation by the angle \(\theta\) defined in Eq.
\(\ref{eq:tan2theta}\).

The off-diagonal matrix element between the dressed states is

\[
\mu_{12}
\equiv
\langle 1|\hat\mu|2\rangle
=
\frac{\mu_e-\mu_g}{2}\sin 2\theta
+
m_0 \cos 2\theta.
\label{eq:mu12general}
\]

Eq. \(\ref{eq:mu12general}\) is the second key result of the note.

It shows that the magnetic transition between the dressed states can
survive in two different ways:

\begin{enumerate}
\def\labelenumi{\arabic{enumi}.}
\tightlist
\item
  \textbf{State mixing}: even if \(m_0 = 0\), one still gets
\end{enumerate}

\[
\mu_{12}
=
\frac{\mu_e-\mu_g}{2}\sin 2\theta.
\label{eq:mu12mixingOnly}
\]

Thus a purely diagonal magnetic operator in the
\(\{|g\rangle,|e\rangle\}\) basis becomes off-diagonal in the dressed
basis whenever the strong Hamiltonian mixes the states.

\begin{enumerate}
\def\labelenumi{\arabic{enumi}.}
\setcounter{enumi}{1}
\tightlist
\item
  \textbf{Direct transition term}: if \(m_0 \neq 0\), then even without
  much mixing one retains the term \(m_0\cos 2\theta\).
\end{enumerate}

The main pedagogical point is therefore:

\begin{quote}
diagonalising the pion-dressed \textbf{Hamiltonian} does not generally
diagonalise the magnetic \textbf{operator}.
\end{quote}

That is why a magnetic transition can remain between the pion-dressed
nuclear states.

\subsection{7) Final post-pion Hamiltonian with the photon still
explicit}\label{final-post-pion-hamiltonian-with-the-photon-still-explicit}

We now keep the photon mode explicit and rewrite the full post-pion
theory in the dressed basis:

\[
H_{\rm post\text{-}\pi}
=
\omega_\gamma a^\dagger a
+
E_1 |1\rangle\langle 1|
+
E_2 |2\rangle\langle 2|
+
g_\gamma (a+a^\dagger)\hat\mu_{\rm dressed}.
\label{eq:postPionHamiltonian}
\]

The transition part of the photon coupling is

\[
H_{\gamma,\rm trans}
=
g_\gamma (a+a^\dagger)\mu_{12}
\Big(
|1\rangle\langle 2|
+
|2\rangle\langle 1|
\Big).
\label{eq:transitionPiece}
\]

So the photon couples states such as

\[
|2,n_\gamma\rangle
\leftrightarrow
|1,n_\gamma+1\rangle
\]

with matrix element

\[
g_\gamma \mu_{12}\sqrt{n_\gamma+1}.
\label{eq:matrixElementPhoton}
\]

Near resonance, the relevant \(2\times 2\) block in the product basis

\[
\{|2,n\rangle,\ |1,n+1\rangle\}
\]

is

\[
\begin{pmatrix}
E_2 + n\omega_\gamma
&
g_\gamma \mu_{12}\sqrt{n+1}
\\
g_\gamma \mu_{12}\sqrt{n+1}
&
E_1 + (n+1)\omega_\gamma
\end{pmatrix}.
\label{eq:photonBlock}
\]

This is the time-independent, quantised-field version of the magnetic
transition. The pions have been integrated out, but the magnetic
coupling between the resulting dressed nuclear states is still present
explicitly through \(\mu_{12}\).

\subsection{8) Special case: a degenerate spin-up / spin-down
doublet}\label{special-case-a-degenerate-spin-up-spin-down-doublet}

For the cleanest version of the toy model, let the retained states be
the same bare nuclear configuration with opposite spin projection:

\[
|g\rangle = |\phi,\downarrow;0_\pi\rangle,
\qquad
|e\rangle = |\phi,\uparrow;0_\pi\rangle.
\label{eq:spinflipstates}
\]

This is the case where it is most natural to take the bare energies
equal:

\[
E_g = E_e \equiv E_0.
\label{eq:spinflipdeg}
\]

The doorway state is still the higher-energy pionful configuration

\[
|\chi,1_\pi\rangle,
\label{eq:spinflipdoorway}
\]

so \(\chi\) labels the nuclear part of the virtual intermediate state,
while the explicit pion sits in the \(1_\pi\) factor.

Because the retained states are degenerate, the denominator in Eq.
\(\ref{eq:BHtoy}\) can be expanded about the common retained-space
energy \(E_0\). Define

\[
\Delta \equiv E_\chi + \omega_\pi - E_0.
\label{eq:spinflipgap}
\]

Then Eq. \(\ref{eq:HeffNuclear}\) simplifies to

\[
H_{\rm nuc}^{\rm eff}
=
E_0
\begin{pmatrix}
1 & 0 \\
0 & 1
\end{pmatrix}
-
\frac{1}{\Delta}
\begin{pmatrix}
W_g^2 & W_gW_e \\
W_gW_e & W_e^2
\end{pmatrix}.
\label{eq:spinflipHeff}
\]

So the pionful sector still does two things:

\begin{itemize}
\tightlist
\item
  it shifts the diagonal energies,
\item
  it mixes the two spin states.
\end{itemize}

Define

\[
\Omega \equiv \sqrt{W_g^2 + W_e^2}.
\label{eq:spinflipOmega}
\]

Then the dressed eigenstates are

\[
|1\rangle
=
\frac{W_g|g\rangle + W_e|e\rangle}{\Omega},
\qquad
|2\rangle
=
\frac{-W_e|g\rangle + W_g|e\rangle}{\Omega},
\label{eq:spinflipdressed}
\]

with energies

\[
E_1 = E_0 - \frac{\Omega^2}{\Delta},
\qquad
E_2 = E_0.
\label{eq:spinflipenergies}
\]

So one linear combination is pulled down by the virtual pionful sector,
while the orthogonal combination is untouched at this order. If \(E_0\)
is interpreted as the free or threshold energy, this gives the simple
toy picture of one bound combination and one state still sitting at
threshold.

Now choose the magnetic operator to be diagonal in the bare spin basis:

\[
\hat{\mu}
=
\begin{pmatrix}
-\mu_0 & 0 \\
0 & \mu_0
\end{pmatrix}_{\{|g\rangle,|e\rangle\}}.
\label{eq:spinflipmu}
\]

This is the practical advantage of the spin-up / spin-down basis: the
two bare states are energy-degenerate, but the magnetic operator still
distinguishes them cleanly.

Transforming Eq. \(\ref{eq:spinflipmu}\) into the dressed basis gives

\[
\mu_{12}
\equiv
\langle 1|\hat{\mu}|2\rangle
=
\mu_0\,\frac{2W_gW_e}{W_g^2+W_e^2}.
\label{eq:spinflipmu12}
\]

So the quantised photon still couples the dressed states as long as both
doorway couplings are nonzero.

\subsection{9) What this toy model does and does not
say}\label{what-this-toy-model-does-and-does-not-say}

What the toy model captures:

\begin{itemize}
\tightlist
\item
  virtual pionful intermediate states can be eliminated,
\item
  their effects survive as level shifts and induced mixing inside a
  no-pion low-energy Hamiltonian,
\item
  the resulting dressed nuclear states can still couple magnetically to
  a quantised photon mode.
\end{itemize}

What the toy model does \textbf{not} yet attempt:

\begin{itemize}
\tightlist
\item
  a realistic two-nucleon spin-isospin structure,
\item
  tensor mixing such as \(^3S_1\)-\(^3D_1\),
\item
  exchange-current corrections derived consistently with a realistic
  nuclear force.
\end{itemize}

Those are natural next extensions, but they are not needed to understand
the basic mechanism shown here.

\subsection{10) One-line summary}\label{one-line-summary}

The full toy theory contains explicit pion and photon modes, but after
integrating out the pionful doorway state one is left with a
\(2\times 2\) pion-dressed nuclear Hamiltonian and a photon coupling
whose off-diagonal matrix element

\[
\mu_{12}
=
\frac{\mu_e-\mu_g}{2}\sin 2\theta
+
m_0\cos 2\theta
\label{eq:finalSummaryFormula}
\]

still drives transitions between the dressed nuclear states.

\printbibliography


\end{document}
